% Einheit 1 von 4 – Stammfunktion (bank/stammfunktion-und-hauptsatz/e1.jsonl, zusammenbau v0.1)
%% TODO Zeitmarke Einheit 1: Marken-Zeile nicht lesbar
%% TODO Zweigzeile Teil 1: was hier gelernt wird (Satz in Schülersprache aus dem Einheitstitel) – Einheit 1
\einheitenkopf[e1]{Einheit 1 von 4 · Stammfunktion}
\zweigzeile{Abitur GK}
\setcounter{aufgabe}{7}

%% TODO Titel: Ich-kann-Satz fehlt in der Bank (Platzhalter: Kettenname) – Nr. 8
%% TODO Anweisung: ein Satz über der ersten Teilaufgabe fehlt in der Bank – Nr. 8
\begin{aufgabe}{Stammfunktion}
%% TODO \rechenplatz{n}: Schrittzahl des Grundfalls fehlt in der Bank (nur bei fester Schreibform, 2.3 b)
\begin{teile}
\teil Gesucht ist eine Funktion $F$ mit $F'(x) = 2x^3 - 5$. Ableiten oder aufleiten? \\ \kreuz{ableiten} \\ \kreuz{aufleiten}
\teil $f(x) = x \cdot (x + 4)$ – eine Stammfunktion $F$? $F(x) = $ \leerfeld
\teil $F(x) = 3x^4 - 2x + 5$ ist eine Stammfunktion von $f$. Gib $f(x)$ an. \hfill (Abitur 2020 GK) \\ $f(x) = $ \leerfeld
\teil Weise nach, dass $F(x) = (-x - 3) \cdot e^{-x}$ eine Stammfunktion von $f(x) = (x + 2) \cdot e^{-x}$ ist, und begründe, dass $f$ nur bei $x = -2$ eine Nullstelle hat. \hfill (Abitur 2018 GK)
\teil $f(x) = 6x^2 - 4x$. Bestimme die Stammfunktion $F$ von $f$, deren Graph durch $P(2 | 5)$ verläuft. \hfill (Abitur 2025 GK) \\ $F(x) = $ \leerfeld
\teil $F(x) = x^2 - 2x$ ist eine Stammfunktion von $f$. Gib alle Stammfunktionen von $f$ an, deren Graph überall den senkrechten Abstand $3$ zum Graphen von $F$ hat.
\teil $f(x) = x \cdot e^{-\frac{x}{2}}$. Für welche Zahl $r$ ist $F(x) = r \cdot (x + 2) \cdot e^{-\frac{x}{2}}$ eine Stammfunktion von $f$? $r = $ \leerfeld
\end{teile}
\end{aufgabe}

%% TODO Titel: Ich-kann-Satz fehlt in der Bank (Platzhalter: Kettenname) – Nr. 9
%% TODO Anweisung: ein Satz über der ersten Teilaufgabe fehlt in der Bank – Nr. 9
\begin{aufgabe}{Stammfunktion – Fehler finden, Begründen}
\begin{teile}
\teil Ich finde den Fehler: $F(x) = x^3 - 4x$ ist eine Stammfunktion von $f$. Tom gibt $f(x) = \frac{1}{4}x^4 - 2x^2$ an.
\teil Begründe, dass $F(x) = x^2 + 5$ und $G(x) = x^2 - 1$ beide Stammfunktionen von $f(x) = 2x$ sind.
\end{teile}
\end{aufgabe}
